\documentclass[a4paper,10pt]{article}
\usepackage{fontspec}
\usepackage{setspace}
\usepackage{biblatex}

\setmainfont{arial}
\setstretch{1.15}
\addbibresource{references.bib}

\title{Artificial Intelligence, Machine Learning and Deep Learning}
\author{Arif Suganda\\
        Tadulako University\\
        \texttt{f5522530002@untad.ac.id}
}

\begin{document}
\maketitle

\section{Artifial Intelligence}

Artificial intelligence (AI) is the capability of computational system to perform task mimicking human intelligence and cognitive function.
These capabilities includes task such as learning, reasoning, and problem solving, which is divided into three main categories of AI, namely:
\begin{enumerate}
    \item Artificial Narrow Intelligence (ANI)
    \item Artificial General Intelligence (AGI)
    \item Artificial Super Intelligence (ASI)
\end{enumerate}

The difference between the three is that ANI has limited capabilities, restricted into one specialty.
The example of ANI is the ASI listed in the cancer detection advacement~\cite{Darbandi2026Cancer}, the scope of this AI limited only for breast cancer.
While AGI and ASI is much more advance than ANI, AGI has the capability of a normal human by the combination of multiple agent for many different cases and ASI is know as super intelligence surpasing any human.
AGI and ASI currently does not exist in any form yet, but the research in this field is still ongoing~\cite{ibm_ai_comparison}.

So how does the computer achieve these capabilities? 
The process of acquiring these knowledge is called machine learning, which revolutionized the way we train the computer today to learn from available data.

\section{Machine Learning}

Machine learning (ML) is a subset artificial intelligence used to train a computer to learn, generalize and acquire intuition from data.
The model created in this training then can be used as a references to help with the decision maker.
There are many technique that can be used by machine learning, such as:
\begin{enumerate}
    \item Linear regression
    \item Bayesian Logic
    \item Support vector machine
    \item Deep learning
\end{enumerate}

Machine learning works by using data as learning source, forming intuition, generalization and adapting to the dataset.
For example the usage of multiple machine learning technique for text classification presented in the survey~\cite{Rady2024}.
In the recent years, machine learning has been heavily dominated by deep learning approach, which inspired by the working of the human brain.

\section{Deep Learning}

Deep learning (DL) is a specialized subset of machine learning, using deeply connected neural network to perform task such as classification, regression and representation learning.
Word ``deep'' refers to the number of layers that the network employs to learn from the dataset, ranging from three to several hundred or thousands.

Deep learning works by clustering the feature found in the data into the interconnected layer and weight in the network.
The clustering process is nearly imposible to interpret or explain, because of this DL often called a blackbox method~\cite{ibm_ai_comparison}.
To understand DL, we can start with understanding the component of neural network, namely:
\begin{enumerate}
    \item Input, used as starting value from the data.
    \item Weight, used to store the adaptive knowledge of the network.
    \item Activation function, used to process the connection of the network and detemined the active connection.
    \item Loss function, used to calculate the rate of error in the network and propagate the value to the weight of the network.
    \item Backpropagation, used to propagate the loss function to adapt the weight and update it into a new value.
\end{enumerate}

Components above working in large network somehow able to adapting to the dataset and found an intricate pattern within the data, then able to adapt and produce an acceptable result.
For example eigen~\cite{eigen14062283} dwelve into depth estimation using deep learning model called convolutional neural network (CNN), although it has an acceptable accuracy, the reseacher still had a hard time to figure how the model works.

\section{Discussion}

The difference between three concept above is in their concept and scope.
Artifial intelligence is an result of learning, machine learning is the theoretical process and deep learning is the learning technique.
Deep learning is a subset of machine learning, machine learning is a subset of artificial intelligence.

\section{Implementation}

AI especially using deep learning can be implemented in a few key problem within the university, namely:
\begin{enumerate}
    \item Attendance, using deep learning to implement face recognition.
    \item Cheater detection, using deep learning to take a realtime video as an input to detect movement that indicate cheating.
    \item AI academic chatbot, using deep learning to create a chatbot to help students or staffs to navigate academic activity in the university.
\end{enumerate}

\printbibliography{}

\end{document}